\section{Conclusion and Discussion}
\label{sec:conclusion}

This work re-examines anisotropy in large language model training
through the lens of activation bias.
While prior analyses focused on spectral spike--tail structure,
we show that a coherent mean component in activations
is the primary driver of extreme elementwise outliers
and second-moment inflation.

Empirically, we demonstrated that:

\begin{itemize}
    \item Activation means are statistically significant across layers and
    training stages, while gradient-output means remain at noise level.
    \item Extreme activation outliers are overwhelmingly dominated by
    the mean component rather than spike variance or semantic tail.
    \item Bias propagates multiplicatively into gradient outer products,
    inflating both blockwise quantization scale and adaptive optimizer
    second moments.
\end{itemize}

These findings clarify that many training instabilities attributed to
spectral anisotropy are, in practice, bias-driven.

\paragraph{Curse.}
Bias introduces dynamic-range inflation in low-bit quantization,
suppresses effective learning rates in adaptive optimizers,
and amplifies extreme coordinate statistics.
Its coherent nature causes disproportionate influence on
blockwise scaling and second-moment estimates.

\paragraph{Blessing.}
Because bias is structured, low-rank, and globally coherent,
it admits an extremely simple mitigation strategy:
mean removal.
Unlike full SVD-based spike suppression,
mean removal is hardware-friendly, inexpensive,
and conservative — it removes the dominant instability
without disturbing semantic covariance structure.

Across large-scale experiments ($1$B parameters, $100$B tokens),
bias mitigation improved both FP4 quantization stability
and adaptive optimizer behavior.
Spike whitening provided further refinement,
but the majority of stabilization arose from bias removal alone.

\subsection{Implications}

The results suggest a revised perspective on LLM training geometry:

\begin{enumerate}
    \item Spectral spike structure is not monolithic.
    A significant portion of spike energy corresponds to bias.
    \item Controlling bias is a minimal structural correction
    that addresses multiple downstream instabilities simultaneously.
    \item Expensive spectral decompositions may be solving
    a simpler underlying structural issue.
\end{enumerate}

This reframing opens new design directions
for quantization schemes, optimizers,
and potentially even architectural normalization strategies.

\subsection{Limitations and Open Questions}

Several important questions remain open:

\begin{itemize}
    \item \textbf{Origin of activation bias.}
    Why do forward activations consistently develop large coherent means?
    Is this purely data-driven, architecture-induced,
    or an emergent property of residual accumulation?
    
    \item \textbf{Interaction with normalization layers.}
    How does bias behave under different normalization schemes
    (pre-norm vs post-norm, RMSNorm vs LayerNorm)?
    
    \item \textbf{Scaling behavior.}
    Does bias magnitude scale with model size or training duration?
    Preliminary evidence suggests early-layer dominance,
    but a systematic scaling law analysis is warranted.
    
    \item \textbf{Beyond mean removal.}
    Can lightweight structured corrections approximate full whitening
    without expensive decompositions?
\end{itemize}

Addressing these questions will further clarify the geometric
role of bias in representation learning.

\subsection{Final Perspective}

Bias is neither purely pathological nor purely benign.
It simultaneously destabilizes training
and provides a coherent structural handle for mitigation.

Recognizing this duality
allows us to move beyond viewing anisotropy
as an opaque spectral phenomenon,
toward understanding and controlling
the specific geometric mechanisms that drive it.

In this sense, bias embodies both
the curse and the blessing
of large language model training.